\section{NLS: the baseline approximation and its instability}
\label{sec:nls}
\subsection{Dysthe's dimensionless variables}
\label{sec:scaling}
From this point, use the transformation in (2.18). To make its meaning
unambiguous, write
\begin{align*}
 t&=\omega_0t_{\rm dim},& (x,y,z)&=k_0(x_{\rm dim},y_{\rm dim},z_{\rm dim}),\\
 A&=\frac{2k_0^2}{\omega_0}A_{\rm dim},&
 \Ph&=\frac{2k_0^2}{\omega_0}\Ph_{\rm dim},&
 \Zh&=k_0\Zh_{\rm dim}.
\end{align*}
The carrier phase becomes $x-t$ and $c_g=1/2$. These are carrier units,
not yet the separate slow coordinates $X,Y,T$. The envelope remains
$O(\varepsilon)$ and varies slowly in $x,y,t$.

Deleting the right side of (2.19) gives
\begin{equation}\label{eq:nls}
 2i\left(A_t+\frac12A_x\right)
 +\frac12A_{yy}-\frac14A_{xx}-|A|^2A=0.
\end{equation}
In the group coordinate $X=x-t/2$, a longitudinal envelope obeys
\[
 iA_t-\frac18A_{XX}-\frac12|A|^2A=0.
\]
Both coefficients are negative. Their product is positive, so this has
the focusing NLS sign combination; reversing time gives the frequently
used convention with two positive signs. A sign difference between
references need not imply different physics, but conventions must be checked.

The dispersive coefficients follow from \eqref{eq:dispersion-series}.
The cubic term describes an amplitude-dependent frequency shift. A
constant envelope solution is
\begin{equation}\label{eq:plane-wave}
 A=A_0e^{-iA_0^2t/2},\qquad A_0>0.
\end{equation}
Its magnitude is constant. Multiplication by the carrier gives total phase
$x-(1+A_0^2/2)t$ to this order. The extra phase rotation is the weakly
nonlinear Stokes frequency correction.

\subsection{What linear stability asks}
Stability asks whether a small disturbance of this particular solution
remains small initially or grows. It does not ask whether $|A|$ itself
oscillates, and it does not prove later breaking.

Introduce an independent small perturbation parameter $\eta$:
\begin{equation}\label{eq:perturbation}
 A=A_0[1+\eta u(x,y,t)]
       e^{i(\eta v(x,y,t)-A_0^2t/2)},
\end{equation}
where $u,v$ are real and $\eta\ll1$. To first order,
$A=A_0e^{-iA_0^2t/2}[1+\eta(u+iv)]+O(\eta^2)$.
The amplitude-expansion parameter $\varepsilon$ and perturbation parameter
$\eta$ have different jobs. Linearization discards $\eta^2$ while retaining
the prescribed powers of $A_0=O(\varepsilon)$.

Set $\Dg=\partial_t+\tfrac12\partial_x$ and
$P_2=\tfrac14\partial_y^2-\tfrac18\partial_x^2$.
Substituting \eqref{eq:perturbation} into \eqref{eq:nls}, cancelling the
background solution, and separating real and imaginary parts gives
\begin{equation}\label{eq:nls-linear}
 \Dg u+P_2v=0,\qquad
 \Dg v-P_2u+A_0^2u=0.
\end{equation}
One useful intermediate identity is
$|A|^2A=A_0^3e^{-iA_0^2t/2}[1+\eta(3u+iv)]+O(\eta^2)$.
The phase shift from the background cancels one part of this contribution.

\subsection{The Fourier-mode matrix calculation}
Use a complex representation for the perturbation modes and take real
parts at the end:
\[
 (u,v)=(\widehat u,\widehat v)e^{i(\lambda x+\mu y-\Omega t)}.
\]
Define
\[
 p=\frac{\lambda^2}{8}-\frac{\mu^2}{4},\qquad
 V=\Omega-\frac\lambda2.
\]
$\Dg$ acts as $-iV$ and $P_2$ acts as $p$. The homogeneous system is
\[
 \begin{pmatrix}-iV&p\\A_0^2-p&-iV\end{pmatrix}
 \begin{pmatrix}\widehat u\\\widehat v\end{pmatrix}=0.
\]
A nonzero perturbation requires determinant zero:
\begin{equation}\label{eq:nls-dispersion}
 \left(\Omega-\frac\lambda2\right)^2=p(p-A_0^2).
\end{equation}

\begin{proposition}[Growth in the linearized NLS model]
For $A_0>0$ and real $\lambda,\mu$, a growing Fourier mode exists precisely
when $0<p<A_0^2$. Its growth rate is
$\gamma=\sqrt{p(A_0^2-p)}$.
\end{proposition}
\begin{proof}
If $0<p<A_0^2$, the right side of \eqref{eq:nls-dispersion} is negative,
so one frequency has $\operatorname{Im}\Omega>0$. For
$\Omega=\Omega_r+i\gamma$, the time factor is
$e^{-i\Omega t}=e^{-i\Omega_rt}e^{\gamma t}$. Otherwise the roots are
real, including repeated roots at the boundaries, and have no positive
exponential growth rate.\qedhere
\end{proof}
Repeated boundary roots can require separate analysis of neutral or
algebraic behavior. The proposition concerns exponential mode growth only.

For longitudinal perturbations, $\mu=0$ and $p=\lambda^2/8$. Thus
\begin{equation}\label{eq:nls-results}
 0<|\lambda|<2\sqrt2 A_0,\qquad
 |\lambda_M|=2A_0,\qquad \gamma_M=\frac12A_0^2.
\end{equation}
To maximize growth, complete the square:
$p(A_0^2-p)=A_0^4/4-(p-A_0^2/2)^2$.
The fastest NLS disturbance therefore has $p=A_0^2/2$.

\subsection{Why transverse perturbations matter}
For two horizontal modulation directions, the maximum condition is
\begin{equation}\label{eq:hyperbola}
 \lambda^2-2\mu^2=4A_0^2.
\end{equation}
This hyperbola is unbounded. The NLS polynomial dispersion approximation
formally predicts equal maximum growth at arbitrarily large modulation
wavenumbers along it, although those wavenumbers violate the narrowband
assumption. Dysthe's correction reduces this degeneracy. It does not make
a small-bandwidth approximation valid at large bandwidth.

\paperlink{Read the introduction's discussion of NLS, then (3.1)--(3.7).
The matrix calculation above is the simpler version of Dysthe's calculation.}
\checkpoint{Explain why $\Omega$ with positive imaginary part grows for
the convention $e^{-i\Omega t}$.}
\clearpage
